\section{Reproducibility and Version Boundaries}
\label{app:execution}

\subsection{Generation Configuration}

All current result rows use the frozen 305-task objects in the same order and
the reference-relative protocol of Section~\ref{sec:judge}. The four native
model groups share the image with SHA-256 prefix
\texttt{1f72e58f6d9b5176} (full digest in the release manifest),
OpenCode 1.18.11, and the education-single-turn agent prompt. Each condition
retains a per-task workspace and logs. Native With exposes the designated
candidate package, not all 14 packages. No rubric or evaluation reference is
given to generation. The package mount and any available supporting files are
part of the release manifest.

The provider aliases are \texttt{glm-5-3-flash-260828},
\texttt{deepseek-v4.1-flash}, \texttt{deepseek-v4-pro-ga-260813}, and
\texttt{doubao-seed-2-1-pro-260915}. GLM Flash, Pro, and Doubao use the Ark
\texttt{/api/v3} protocol; DeepSeek Flash uses a third-party OpenAI-compatible
\texttt{/v1} endpoint. GLM Flash enables low-effort thinking, whereas the other
three disable it. These recorded switches are not claims about hidden provider
computation. Environment concurrency is 8 for GLM Flash and DeepSeek Flash and
32 for Pro and Doubao; the API concurrency ceiling is 64. Scheduling
concurrency is not repeated sampling or a per-task roll-out count.

The retained GLM-5.3 row has 262 With answers from full-text Skill exposure
with native resources and 43 from native candidate-package exposure. It is a
historical composite result under the current task inventory and score metric,
not a fifth matched native-exposure intervention. It is retained for result
coverage but excluded from native-access statistics. The task-level exposure
manifest is necessary to reproduce this row; an aggregate score cannot identify
the effect of a single exposure method.

\subsection{Judge Inputs and Aggregation}

The accepted judge uses \texttt{deepseek-v4.1-flash}, thinking disabled,
6,000 maximum output tokens, and a final recovery timeout of 600 seconds.
The fixed reference answer text is the archived Pro No Skill answer for each
task, not a human gold answer. Candidate and reference are anonymized jointly;
the rubric scope, segmentation, A/B ordering, and hashes are retained.
High-scoring evidence citations refer to the appropriate response's own
segments. Score validation checks structure and bindings, not the truth of the
quoted claims. Each accepted pair retains candidate and reference dimension
scores, evidence, and the derived ratio.

The primary calculation is identical for all 305 tasks: equally average
applicable 0--10 dimensions within each response, compute its per-task
candidate/reference ratio, and average those ratios over the full inventory.
An execution failure contributes zero; a reference-self cell contributes 0.5.
Missing or invalid verdicts must remain unresolved, not be assigned convenient
defaults. None remain in the five-model archive. Transport retries and rejected
schema responses may precede an accepted verdict, so one retained verdict does
not imply exactly one historical HTTP request.

\subsection{Historical Rubric Metadata}

The source inventory contains 60 tasks with three native grades, 100 with four,
20 with two grades plus a below-minimum fallback, and 83 with satisfaction
labels where native grades are absent. The remaining 42 constructed tasks have
weighted binary items. This provenance is retained, but the current 305-task
metric does not count top grades on one subset and all-pass outcomes on another.
The current human-adjudicated comparison in Appendix~\ref{app:human} uses
the same continuous dimensions and paired answers as the deployed Ref judge,
not those earlier aggregation formulas. A pointwise score and a normalized
relative score are on different scales; their numerical difference does not
measure a quality change or a correction of inflation.

\subsection{Small-Model Provenance}

Training seeds use the EduBench snapshot at commit
\texttt{f19004a5c68c27cd9d35\allowbreak 4401654345072fb712c6}. This is a training-seed
source, not a claim that every evaluation task is an official EduBench item.
The 246-example release contains 80 Skill--seed combinations from 70 upstream
seeds. Variant counts are 62 complete-product tasks, 65 changed-evidence tasks,
61 resource-constraint tasks, and 58 artifact-repair tasks. The recorded
\texttt{environment\_trace\_verified} flag is false for all 246 examples;
the retained targets are complete final answers, not verified full trajectories.

The SFT base is Qwen3-4B-Instruct-2507. Its one-epoch LoRA run uses bf16 and
completion-only labels, rank 32, alpha 64, dropout 0.05, learning rate $10^{-5}$,
and cosine scheduling with 5\% warmup. Adapters update attention $q/k/v/o$
and MLP gate/up/down projections. Micro-batch 8 with accumulation 2 gives an
effective batch of 16; training completes 16 optimizer steps.
The longest retained training sequence has 7,723 tokens,
within the 8,192-token limit; the mean assistant target has 1,091.81 tokens.
The run takes approximately 239 seconds. GRPO completes 123 optimizer steps
at learning rate $2\times10^{-6}$, using eight samples per question and
generation batch 16, with a 4,096-token output limit, temperature 0.7,
top-$p$ 0.9, cosine scheduling, 5\% warmup, and a DAPO-style loss. Its frozen
KL policy is the SFT adapter with coefficient 0.04; empty and truncated outputs
receive zero reward. Its run takes approximately 3,536 seconds.
The training outcome judge is \texttt{deepseek-v4-flash-ga-260731}; the
evaluation judge is V4.1 Flash. These are different model versions and reward
scales, as specified in Section~\ref{sec:internalization}.

The current review draft's small-model table is bound to the supplied recovery
summary, while the latest item-level recovery archive is pending local
reconciliation. That table should not be finalized for submission until the
answer hashes, accepted verdicts, failure states, and paired differences are
checked. The five API-system tables and native-access tables have already been
recomputed from their retained item-level materials.

\subsection{Release Contents}

The retained experiment archive includes task/source/Skill IDs, exposure modes,
model request aliases, generation and judging settings, accepted responses,
failure states, rubric scopes, joint judgments, access evidence tiers, and
offline aggregation scripts. The manuscript bundle contains compact frozen
tables, case exports, and provenance manifests. These materials separate original pointwise experiments from
current reference-relative experiments. Credentials, personal server routes,
and unrelated projects are excluded. Availability of a record is not silently
promoted to evidence of successful file delivery or scientific correctness.
